\section{GPU Reference-Orbit Implementation (Fused v2)}
\label{subsec:ref-orbit-gpu}

Deep Mandelbrot views require a high-precision reference orbit before the
perturbation renderer can evaluate nearby pixels.  The recurrence
\begin{align}
  x_{n+1} &= x_n^2-y_n^2+c_x, \\
  y_{n+1} &= 2x_ny_n+c_y
  \label{eq:gpu-reference-components}
\end{align}
uses $x_n$ and $y_n$ for the components of the reference orbit at
$c_\star=c_x+i c_y$.
It is simple algebraically, but its cost changes character when each component has
thousands or millions of bits.  Ordinary machine arithmetic is then replaced
by multiplication of long coefficient sequences, followed by carry
propagation and normalization.  At the deepest supported precisions, this work
dominates reference construction.

The production GPU backend implements the complete recurrence inside one
persistent cooperative CUDA kernel.  Packing, number-theoretic transforms
(NTTs), pointwise products, signed carry propagation, normalization, derivative
updates, periodicity checks, and bounded orbit publication are phases of that
engine rather than independently launched arithmetic operators.  This chapter
develops the mathematical motivation for the NTT and then relates it to the
specific fused schedule.

The implementation described here is the only supported GPU backend.  It is
called the second-generation implementation, Ref2, or fused v2.  The removed v1
backend composed separate high-precision multiply and add/subtract operators.
Unless a test symbol is quoted explicitly, \emph{GPU reference orbit} hereafter
means the fused v2 engine.

\subsection{The Parallelism Available in a Sequential Orbit}
\label{subsec:gpu-reference-parallelism}

The orbit itself cannot be divided among iterations: iteration~$n+1$ needs the
normalized value produced by iteration~$n$.  Assigning one CUDA thread to each
iteration would therefore violate the data dependency rather than expose useful
parallelism.  The opportunity lies \emph{within} one iteration.  A long
significand contains many limbs, a multiplication contains many coefficient
products, and an NTT contains many independent butterflies at every stage.
The kernel advances one orbit in time while the cooperative grid evaluates each
high-precision step in space.

This decomposition creates five design requirements:

\begin{enumerate}
  \item Products must be exact at the selected finite precision.  An unnoticed
        error can be amplified over a long orbit or corrupt a Newton step.
  \item The arithmetic should present many uniform operations to the GPU rather
        than a small number of recursive, irregular tasks.
  \item Intermediate products should remain in a reusable workspace instead of
        becoming separately normalized high-precision values.
  \item The unavoidable transform and carry synchronization should be shared
        across the real, imaginary, and optional derivative streams.
  \item The completed value must return to positional form after every
        iteration because the next plan depends on its exponent, sign, and
        normalized significand.
\end{enumerate}

\subsubsection{Why transform multiplication is used}

Several multiplication families can satisfy some of these requirements, but
their GPU behavior differs.  \Cref{tab:gpu-multiply-choices} summarizes the
design tradeoff.  The asymptotic costs describe multiplication of coefficient
sequences of length~$L$; they do not imply that the NTT is fastest for every
small precision.

\begin{table}[!htbp]
\centering
\small
\caption[Multiplication choices for a GPU reference orbit]{Multiplication choices for a GPU
  reference orbit.  The production path favors exactness and regular parallel work over the
  smallest setup cost.}
\label{tab:gpu-multiply-choices}
\begin{tabularx}{\textwidth}{lXXX}
\hline
\textbf{Method} & \textbf{Work structure} & \textbf{Advantage} & \textbf{Limitation here} \\
\hline
Schoolbook
  & $O(L^2)$ independent digit products followed by carries
  & Simple and effective for short operands
  & Quadratic work becomes prohibitive as precision grows \\
Karatsuba/Toom
  & Recursive subproblems and recombination
  & Fewer scalar products than schoolbook multiplication
  & Uneven subproblem sizes, extra temporaries, and recursive synchronization complicate a
    persistent SIMT schedule \\
Floating-point FFT
  & Regular $O(L\log L)$ butterflies and pointwise products
  & Mature high-throughput transform techniques
  & Rounded coefficients require a separate error bound and recovery strategy over long orbits \\
Number-theoretic transform
  & Regular $O(L\log L)$ modular butterflies and pointwise products
  & Exact convolution with integer residues and deterministic recovery
  & Modular multiplication, root tables, padding, and stage barriers have nontrivial cost \\
\hline
\end{tabularx}
\end{table}

Karatsuba and related factorizations remain important CPU multiplication
techniques, and floating-point FFTs can map extremely well to GPUs.  The fused
engine nevertheless gives priority to exact coefficient recovery.  An NTT has
the same convolution structure as an FFT, but it operates in a finite field:
there is no transform roundoff to accumulate or to bound.  Its radix-2 stages
also consist of the same additions, subtractions, twiddle multiplications, and
address pattern for a large collection of coefficients.  Those properties make
the algorithm a natural match for SIMT execution.  GPU multi-precision design
tradeoffs are discussed more generally by Emmart~\cite{Emmart_GPU_MultiPrecision}.

\subsection{From Positional Multiplication to Convolution}
\label{subsec:ntt-convolution-introduction}

Let $B=2^b$ and write two nonnegative significands as
\begin{align}
  A &= \sum_{j=0}^{L-1} a_j B^j, &
  C &= \sum_{j=0}^{L-1} c_j B^j,
  \qquad 0\le a_j,c_j < B.
  \label{eq:ntt-positional-inputs}
\end{align}
Their product before carrying is
\begin{equation}
  AC = \sum_{k=0}^{2L-2} h_k B^k,
  \qquad
  h_k = \sum_{\substack{i+j=k\\0\le i,j<L}} a_i c_j.
  \label{eq:ntt-linear-convolution}
\end{equation}
The sequence $h$ is the \emph{linear convolution} of the coefficient sequences
$a$ and~$c$.  Positional multiplication and polynomial multiplication are thus
the same calculation until the coefficients are carried back into the range
$[0,B)$.

A small decimal example illustrates the separation.  With $B=10$,
\begin{equation}
  (3+5B+2B^2)(4+B)=12+23B+13B^2+2B^3.
\end{equation}
The convolution coefficients $(12,23,13,2)$ are not decimal digits.  Carrying
from low to high gives $12=2+1B$, then $23+1=4+2B$, then
$13+2=5+1B$, and finally $2+1=3$.  The resulting digits are
$(2,4,5,3)$, and hence
$253\mathbin{\times}14=3542$.  The GPU engine follows the same two-part
strategy at much larger scale: the NTT computes all pre-carry coefficients in
parallel, and a prefix algorithm subsequently resolves the positional carry.

Direct evaluation of~\eqref{eq:ntt-linear-convolution} performs $O(L^2)$
coefficient products.  Transform multiplication instead evaluates the two
polynomials at a carefully selected set of field points, multiplies the values
pointwise, and interpolates.  The forward and inverse transforms require
$O(M\log M)$ butterflies for a padded transform length~$M$, while the pointwise
stage requires $O(M)$ products.  Pollard's finite-field formulation gives the
classical basis for this construction~\cite{Pollard1971FiniteFieldFFT}.

\subsection{Number-Theoretic Transform Fundamentals}
\label{sec:ntt-multiply}

Let $p$ be prime, let $M$ be a power of two dividing $p-1$, and let
$\omega\in\mathbb{F}_p$ have order exactly~$M$.  Thus
\begin{equation}
  \omega^M=1\pmod p,
  \qquad
  \omega^j\ne 1\pmod p\quad(0<j<M).
  \label{eq:ntt-primitive-root}
\end{equation}
After zero-padding $a$ to length~$M$, its forward NTT is
\begin{equation}
  \widehat a_k = \sum_{j=0}^{M-1} a_j\omega^{jk}\pmod p,
  \qquad 0\le k<M.
  \label{eq:ntt-forward-definition}
\end{equation}
The inverse transform uses $\omega^{-1}$ and the modular inverse of~$M$:
\begin{equation}
  a_j = M^{-1}\sum_{k=0}^{M-1}\widehat a_k\omega^{-jk}\pmod p.
  \label{eq:ntt-inverse-definition}
\end{equation}
These equations are the discrete Fourier transform equations with complex roots
of unity replaced by roots in a finite field.  Every operation is integer
arithmetic modulo~$p$.

The convolution theorem converts~\eqref{eq:ntt-linear-convolution} into
independent pointwise products:
\begin{equation}
  h = \operatorname{NTT}^{-1}
      \bigl(\operatorname{NTT}(a)\odot\operatorname{NTT}(c)\bigr),
  \label{eq:ntt-convolution-theorem}
\end{equation}
where $\odot$ denotes componentwise multiplication.  Equation
~\eqref{eq:ntt-convolution-theorem} naturally computes convolution modulo
$x^M-1$, so a coefficient beyond position $M-1$ would wrap to the beginning.
For unshifted length-$L$ inputs, $M\ge 2L-1$ prevents that aliasing.  The fused
planner uses the stronger condition that $M$ contain every coefficient after
the operands' exponent-alignment shifts have been included.  If the last
nonzero input coefficients have indices $k_a$ and~$k_c$, the last possible
product index is $k_a+k_c$, so the planner requires $M\ge k_a+k_c+1$.
For example, two unshifted length-three inputs can have a degree-four product:
an $M=4$ transform would fold its last coefficient into index zero, whereas
the next power-of-two length, $M=8$, preserves the linear convolution.

\subsubsection{Radix-2 butterflies}

A direct evaluation of~\eqref{eq:ntt-forward-definition} is still quadratic.
The radix-2 fast transform factors it into $\log_2 M$ stages.  A decimation-in-
frequency (DIF) butterfly combines two values $u$ and~$v$ as
\begin{equation}
  u' = u+v\pmod p,
  \qquad
  v' = (u-v)\tau\pmod p,
  \label{eq:ntt-dif-butterfly}
\end{equation}
where the stage and butterfly position select the twiddle~$\tau$.  Each stage
contains $M/2$ independent butterflies.  The inverse decimation-in-time (DIT)
butterfly reverses this factorization with inverse twiddles and the scale in
~\eqref{eq:ntt-inverse-definition}.

\begin{figure}[!htbp]
\centering
\begin{tikzpicture}[
    >=Stealth,
    val/.style={circle, draw, minimum size=0.62cm, inner sep=1pt, font=\small},
    op/.style={rounded corners, draw, fill=blue!8, minimum width=2.4cm,
              minimum height=0.72cm, align=center, font=\small},
]
  \node[val] (u) at (0, 0.75) {$u$};
  \node[val] (v) at (0, -0.75) {$v$};
  \node[op] (add) at (3.0, 0.75) {$u+v\pmod p$};
  \node[op] (sub) at (3.0, -0.75) {$(u-v)\tau\pmod p$};
  \node[val] (uo) at (5.5, 0.75) {$u'$};
  \node[val] (vo) at (5.5, -0.75) {$v'$};

  \draw[->] (u) -- (add);
  \draw[->] (v) -- (add);
  \draw[->] (u) -- (sub);
  \draw[->] (v) -- (sub);
  \draw[->] (add) -- (uo);
  \draw[->] (sub) -- (vo);
\end{tikzpicture}
\caption[Radix-2 NTT butterfly]{One forward DIF NTT butterfly.  A stage applies this fixed
  operation to $M/2$ independent input pairs, with position-dependent twiddles~$\tau$.}
\label{fig:ntt-butterfly}
\end{figure}

The regularity is as important as the asymptotic reduction.  Within a stage,
threads execute the same modular operations and differ principally in their
coefficient indices and twiddle values.  There is no data-dependent recursion,
and the power-of-two index relationships can be generated with shifts and
masks.

\subsubsection{Exact recovery from residues}

The inverse transform returns each convolution coefficient modulo~$p$.  A
residue $r$ has a unique centered interpretation
\begin{equation}
  \operatorname{center}(r)=
  \begin{cases}
    r, & 0\le r < p/2, \\
    r-p, & p/2 < r < p,
  \end{cases}
  \label{eq:ntt-centered-residue}
\end{equation}
provided that the true signed coefficient has absolute value below~$p/2$.
The base, transform ceiling, and fused product formulas are selected so that
the worst supported coefficient sum retains this headroom.  Consequently,
centered conversion recovers an integer coefficient, not an approximation to
one.  This exactness is the central reason for choosing an NTT over a
floating-point FFT in the reference engine.

\subsection{Why the NTT Maps Well to CUDA}
\label{subsec:ntt-gpu-mapping}

CUDA executes groups of threads in a SIMT hierarchy: lanes form warps, warps
form thread blocks, and a cooperative launch can synchronize the complete grid.
The official CUDA programming model describes both these group levels and the
residency requirement for grid-wide synchronization~\cite{NvidiaCudaProgrammingGuide}.
The NTT supplies useful work at each level.

\paragraph{Grid-level parallelism.}
Large butterfly stages address pairs that may be far apart.  The spectra reside
in global memory, and the cooperative grid distributes the transform work.
The release matrix schedule groups several butterfly stages into each phase;
grid-wide barriers separate phases that exchange data between blocks.  Local
stages reuse shared memory and warp registers.  The memory pattern remains
regular even when butterfly partners are distant.

\paragraph{Block-level locality.}
The release matrix schedule limits each shared-memory dimension to 2,048
coefficients in TwoWay mode.  FourWay mode permits up to 1,024 with 48~KiB of
dynamic shared memory, or 2,048 with at least 96~KiB.  The transform stages are
divided approximately evenly among dimensions, so actual dimensions can be
smaller: a $2^{16}$ transform uses two 256-coefficient dimensions.
One thread block stages the active streams for a local dimension.  Several
stages can then reuse those coefficients and cached twiddles without a
global-memory round trip.  Block barriers replace grid barriers inside the tile.  The same
shared arena is reused by forward, pointwise, inverse, and carry phases, so its
allocation cost is paid once per resident block rather than once per operator.

\paragraph{Warp-level exchange.}
When a complete 64-coefficient batch is available, the innermost five forward
stages, the pointwise products, and the first five inverse stages remain in
warp registers.  Shuffle operations exchange butterfly partners without a
block barrier.  Tiles that do not satisfy the eligibility conditions retain
the shared-memory path, preserving correctness for short or partial tails.

\paragraph{Concurrent streams.}
The orbit-only path transforms $x$ and $y$ in TwoWay mode.  Newton--Raphson
evaluation transforms $x$, $y$, $dx$, and $dy$ in FourWay mode.  All enabled
streams have the same $M$, roots, stage order, and dependency boundaries.
One thread can therefore apply a loaded twiddle to several streams before
moving to the next coefficient, and the additional derivative work does not
require a separate kernel schedule.

The match is not perfect.  One reference orbit supplies only one transform at
a time, so available grid parallelism is bounded by~$M$.  Shared memory and
registers constrain occupancy, modular multiplication has a dependency chain,
and large stages still require device-wide barriers.  The fused design is
therefore not based on an assumption that synchronization is free.  It reduces
the number of separately synchronized operators and places data at the
narrowest synchronization level that can legally contain each stage.

\subsection{High-Precision Representation and Alignment}
\label{subsec:hpfloat-model}

An \code{HpSharkFloat} value consists of a sign, a base-$2^{32}$ significand,
and an exponent.  If the significand limbs encode the nonnegative integer~$A$
and the sign bit encodes $s\in\{-1,1\}$, the represented value is
\begin{equation}
  v=sA2^e.
  \label{eq:hp-shark-value}
\end{equation}
The generated type fixes a compile-time storage bucket
\code{GlobalNumUint32}.  A runtime request selects an active limb count~$D$
greater than half the bucket size and at most its full size.  Arithmetic reads
the most significant $D$ limbs and ignores lower limbs, allowing one generated
specialization to serve a range of nearby precisions.

The NTT plan splits each active 32-bit limb into two base-$2^{16}$
coefficients.  Thus
\begin{equation}
  b=16,
  \qquad
  B=2^{16},
  \qquad
  L=2D.
  \label{eq:reference-b16-plan}
\end{equation}
The smaller coefficient base increases the number of transform elements but
leaves ample field headroom for the sum of many products and for the signs and
small integer factors in the fused recurrence.

Products and constants generally have different binary exponents.  The grid
leader therefore constructs a new \code{HpSharkReferenceIterationPlan} from the
current state before every attempted recurrence.  For each group of operands, it
chooses the least exponent among the nonzero values as the alignment origin.
An exponent difference is decomposed into a whole coefficient shift and a
residual shift:
\begin{equation}
  \Delta e = qb+r,
  \qquad
  q=\left\lfloor\frac{\Delta e}{b}\right\rfloor,
  \qquad
  0\le r<b.
  \label{eq:reference-alignment-split}
\end{equation}
The whole shift moves the coefficient support by~$q$ positions.  The residual
shift combines neighboring 16-bit halves while packing.  Alignment is thus
performed once in the input representation; the pointwise loop does not need
a frequency-dependent spectral shift.

For a nonzero component, the shift in~\eqref{eq:reference-alignment-split}
gives a last possible coefficient index
\begin{equation}
  k_{\max}=
  \begin{cases}
    q+L-1, & r=0, \\
    q+L, & r>0.
  \end{cases}
  \label{eq:reference-packed-support}
\end{equation}
The extra position in the second case holds bits carried out of the highest
16-bit coefficient by the residual shift.  Thus a real product involving both
$X$ and~$Y$ needs up to $2\max(k_X,k_Y)+1$ coefficients, while $X\star Y$
needs up to $k_X+k_Y+1$.  The same support calculation is applied to the
enabled derivative products.

The planner computes the last potentially nonzero coefficient for every
required product, including the derivative products, and rounds the largest
linear-convolution support up to a power of two.  The active length~$M$ is the
larger of that rounded length and the prepared minimum transform length, and
selects a cached plan.  The initial zero state still receives a plan.  A
zero state with no linear contribution takes the
\code{Zero} path, a state with only $c$ or the derivative's $+1$ takes the
\code{LinearOnly} path, and all nonlinear cases take the \code{Ntt} path.  A
required length beyond $2^{25}$ cannot be represented by the workspace; the
kernel reports an unknown status instead of allowing cyclic wraparound or an
out-of-bounds workspace access.

\subsection{Finite-Field and Transform Setup}
\label{subsec:reference-ntt-field}

The reference engine uses the Goldilocks prime
\begin{equation}
  p=2^{64}-2^{32}+1.
  \label{eq:reference-goldilocks-prime}
\end{equation}
Because
\begin{equation}
  p-1=2^{32}(2^{32}-1),
\end{equation}
the field contains power-of-two roots through order~$2^{32}$, comfortably
beyond the engine's $2^{25}$ workspace ceiling.  The setup kernel derives the
forward and inverse stage roots, fills the twiddle tables, computes the inverse
length, and marks each generated plan valid.  \code{ReferencePreparedTables}
owns these tables and the spectra, limb, magnitude, and carry workspaces.  The
tables are reused across all iterations and chunk invocations of a session.

Field products use Montgomery multiplication~\cite{Montgomery1985Modular}.
Write $R_{\mathrm M}=2^{64}$ for the Montgomery radix, distinct from the real
product $R$ below.  A field element $a$ in Montgomery form is
$\widetilde a=aR_{\mathrm M}\bmod p$, and a Montgomery product computes
\begin{equation}
  \operatorname{Mont}(\widetilde a,\widetilde c)
  =\widetilde a\widetilde c R_{\mathrm M}^{-1}
  =acR_{\mathrm M}\pmod p
\end{equation}
without integer division.  Additions and subtractions remain ordinary modular
operations.

The release path also absorbs transform normalization into packing.  For a
length-$M$ product, the prepared field value~$s$ satisfies
\begin{equation}
  s^2=(MR_{\mathrm M})^{-1}\pmod p.
  \label{eq:reference-input-scale}
\end{equation}
Packing actually multiplies each input coefficient by $sR_{\mathrm M}$,
represented in the code as \code{InputScaleR}.  If $\widehat a_k$ and
$\widehat c_k$ are the transforms of the unscaled inputs, linearity and one
pointwise Montgomery product give
\begin{align}
  \operatorname{NTT}(sR_{\mathrm M}a)_k
    &=sR_{\mathrm M}\widehat a_k, \\
  \operatorname{Mont}(sR_{\mathrm M}\widehat a_k,
                      sR_{\mathrm M}\widehat c_k)
    &=s^2R_{\mathrm M}\widehat a_k\widehat c_k.
\end{align}
The inverse stages omit their usual factor~$M^{-1}$, so the final coefficient
is $Ms^2R_{\mathrm M}h_j=h_j\pmod p$ by
\eqref{eq:reference-input-scale}.  It is already a standard residue ready for
centered conversion.  For an even number $t=\log_2M$ of stages,
$sR_{\mathrm M}=2^{32-t/2}\pmod p$, and packing uses a shift.  Odd stage
counts use a specialized multiply of a 16-bit coefficient by the prepared
constant.  No separate inverse-scale or Montgomery-exit pass is needed.

\subsection{The Fused Reference Recurrence}
\label{subsec:fused-ntt-reference-step}

Let $X$ and $Y$ denote the aligned coefficient polynomials for $x_n$ and
$y_n$.  Although~\eqref{eq:gpu-reference-components} is commonly written as
two squares and one product, the production pointwise stage uses the factored
real component
\begin{align}
  R &= (X+Y)\star(X-Y),
       \label{eq:fused-reference-real} \\
  I &= 2(X\star Y),
       \label{eq:fused-reference-imag}
\end{align}
where $\star$ denotes linear convolution.  These expressions equal
$X\star X-Y\star Y$ and $2X\star Y$ exactly but require only two nonlinear
products.  In the transform domain they become
\begin{align}
  \widehat R_k &= (\widehat X_k+\widehat Y_k)
                  (\widehat X_k-\widehat Y_k), \\
  \widehat I_k &= 2\widehat X_k\widehat Y_k.
  \label{eq:fused-reference-pointwise}
\end{align}
The signs of $x_n$ and $y_n$ have already been incorporated during packing.

The constants $c_x$ and $c_y$ are not transformed.  After the inverse NTT,
the signed gather inserts their significand bits at the offsets determined by
the iteration plan.  Interpreting $R$ and~$I$ at their planned product
exponents, carry propagation therefore implements the complete recurrence
\begin{align}
  x_{n+1} &= (x_n+y_n)(x_n-y_n)+c_x
           =x_n^2-y_n^2+c_x, \\
  y_{n+1} &= 2x_ny_n+c_y
\end{align}
without materializing normalized values for $R$ and~$I$ or launching a
standalone high-precision addition.

\subsubsection{Derivative products}
\label{subsec:fused-ntt-derivatives}

Newton--Raphson support also advances
\begin{equation}
  d_{n+1}=2z_nd_n+1,
  \qquad d_n=d_{x,n}+i d_{y,n}.
  \label{eq:fused-dzdc-complex}
\end{equation}
FourWay mode transforms the aligned derivative coefficient polynomials $D_x$
and $D_y$ with $X$ and~$Y$.  The pointwise center forms three products
\begin{align}
  P_1 &= X\star D_x, &
  P_2 &= Y\star D_y, &
  P_3 &= (X+Y)\star(D_x+D_y),
  \label{eq:fused-derivative-products}
\end{align}
The third product expands to
\begin{equation}
  P_3=P_1+X\star D_y+Y\star D_x+P_2,
\end{equation}
so subtracting $P_1$ and~$P_2$ isolates the two cross terms.  The nonlinear
coefficient polynomials are
\begin{align}
  2(P_1-P_2) &= 2(X\star D_x-Y\star D_y), \\
  2(P_3-P_1-P_2) &= 2(X\star D_y+Y\star D_x).
\end{align}
Interpreting these polynomials at their planned product exponent and inserting
the aligned constant gives the scalar recurrence
\begin{align}
  d_{x,n+1} &= 2(x_n d_{x,n}-y_n d_{y,n})+1, \\
  d_{y,n+1} &= 2(x_n d_{y,n}+y_n d_{x,n}).
  \label{eq:fused-derivative-reconstruction}
\end{align}
The derivative's $+1$ enters during signed gathering just as $c_x$ and~$c_y$
do; it is not an unshifted constant coefficient in the product polynomial.
The Gauss-style product identity reduces four real products to three and
shares every transform boundary with the orbit streams.

The second derivative
\begin{equation}
  q_{n+1}=2(d_n^2+z_nq_n)
\end{equation}
is used by the Feature Finder but does not join the high-precision FourWay NTT.
A leader thread updates its configured HDR representation from truncated orbit
and first-derivative values before the fused high-precision step.  The complete
Newton/Halley use of these values is described in
\cref{subsec:ff-inner-gpu}.

\subsubsection{Operation schedule}
\label{subsec:fused-ntt-operation-schedule}

\Cref{fig:ntt-pipeline} summarizes the release matrix schedule, selected for
$M>1024$ when checksum debugging is disabled.  Smaller transforms and checksum
debugging use the direct DIF/DIT schedule with tiles of at most 1,024
coefficients.  Matrix phases transform local dimensions and transpose data
between global buffers.  The center fuses its forward stages, pointwise
products, and inverse stages, keeping eligible innermost stages in warp
registers.  Red dashed boundaries indicate dependencies between grid-wide
phases; local stages synchronize only within the owning block or warp.

\begin{figure}[!htbp]
\centering
\begin{tikzpicture}[
    >=Stealth,
    phase/.style={draw, rounded corners, text width=2.25cm, inner xsep=0.07cm,
                  minimum height=1.55cm, align=center, font=\scriptsize},
    global/.style={phase, fill=blue!10},
    local/.style={phase, fill=green!12},
    carry/.style={phase, fill=orange!14},
    barrier/.style={draw, red!70!black, line width=1.4pt, dashed},
    barrlbl/.style={font=\tiny\bfseries, red!70!black},
]
  \node[global] (pack) at (0, 0)
    {\textbf{Align and pack}\\TwoWay: $x,y$\\FourWay: $x,y,d_x,d_y$};
  \node[global] (forward) at (2.75, 0)
    {\textbf{Forward phases}\\local DIF\\transpose / store};
  \node[local] (center) at (5.5, 0)
    {\textbf{Fused center}\\shared tiles\\warp-register core};
  \node[global] (inverse) at (8.25, 0)
    {\textbf{Inverse phases}\\local DIT\\transpose / store};
  \node[carry] (finish) at (11.0, 0)
    {\textbf{Signed gather}\\prefix carry\\normalize\\publish};

  \foreach \i/\j in {pack/forward, forward/center, center/inverse, inverse/finish}
    \draw[->] (\i) -- (\j);

  \foreach \x in {4.125, 6.875, 9.625} {
    \draw[barrier] (\x, -0.93) -- (\x, 0.93);
    \node[barrlbl] at (\x, -1.13) {grid sync};
  }

  \node[font=\tiny, green!40!black, align=center] at (5.5, 1.12)
    {block barriers; warp shuffles where eligible};
\end{tikzpicture}
\caption[GPU fused NTT schedule and synchronization hierarchy]{GPU fused NTT schedule for one
  iteration, shown as conceptual phases rather than a fixed barrier count.
  Packing is fused into the first active transform phase.  Forward and inverse
  phases outside the center may repeat or be absent, depending on the dimension
  count.  FourWay mode widens the streams and products and can change that count.}
\label{fig:ntt-pipeline}
\end{figure}

In the matrix path, one nonlinear iteration proceeds as follows:

\begin{enumerate}
  \item The grid leader inspects the normalized state, builds the alignment and
        support plan, and publishes its kind, active~$M$, offsets, and flags.
  \item Threads pack signed base-$2^{16}$ coefficients at their aligned
        positions and apply the input scale from
        ~\eqref{eq:reference-input-scale}.  Packing is fused into the first
        forward matrix phase, or into the center for a single-dimension plan.
  \item Forward matrix phases, when needed, transform local dimensions and
        transpose their outputs for the next phase, with grid synchronization
        between phases.
  \item Each block loads a shared-memory center dimension, completes the small
        forward stages, evaluates~\eqref{eq:fused-reference-pointwise} and, when enabled,
        ~\eqref{eq:fused-derivative-products}, and begins the inverse DIT.
  \item Inverse matrix phases, when needed, reverse the dimension traversal.
        The completed inverse returns natural-order coefficients as standard
        residues in global memory.
  \item The signed gather centers residues, incorporates $c_x$, $c_y$, and the
        derivative's $+1$, and constructs carry transfer functions.
  \item A packed prefix operation resolves the carry chains.  Magnitude
        extraction and normalization publish the next orbit and derivative
        states.
  \item If another iteration is requested, the next loop entry records the
        current low-precision orbit value and checks periodicity and escape
        when enabled.  A bounded batch is returned to the host after the launch.
\end{enumerate}

\subsubsection{DIF/DIT ordering and implicit permutation}

In the direct radix-2 schedule, the forward DIF transform consumes natural-order
coefficients and leaves its spectrum in bit-reversed order.  Pointwise products do not care about that
permutation as long as every stream uses the same order.  The inverse DIT
transform consumes the bit-reversed spectra and emits natural-order
coefficients.  Pairing the two variants therefore avoids a separate bit-reversal
kernel and its global reads, writes, and synchronization.  The matrix schedule
also transposes between dimensions, but applies matching permutations to every
stream and pairs its forward and inverse layouts.  Avoiding a standalone
bit-reversal pass does not eliminate those transposes.

The fused center exploits the same pairing more aggressively.  It completes
the small forward stages, evaluates the pointwise recurrence, and starts the
small inverse stages while a tile is resident.  The data need not be written to
a standalone \emph{forward result} buffer and then loaded by an unrelated
multiply operator.  This is the central meaning of \emph{fused} in v2.

\subsection{Signed Carry, Normalization, and Exactness}
\label{subsec:fused-ntt-carry}
\label{sec:hp-add}

Three independent conditions make the modular result an exact positional
result:

\begin{enumerate}
  \item The active~$M$ contains the complete shifted linear-convolution support,
        so no nonzero coefficient wraps around $x^M-1$.
  \item The absolute coefficient bound for every enabled orbit and derivative
        combination is below $p/2$, so
        ~\eqref{eq:ntt-centered-residue} recovers the unique signed integer.
  \item Carry propagation includes every nonlinear and aligned linear
        contribution before normalization truncates the result to its storage
        bucket; subsequent arithmetic reads the selected $D$ most significant limbs.
\end{enumerate}

The coefficient bound can be made explicit for the nonlinear terms.  Each
unscaled aligned input coefficient has absolute value at most $B-1$, and each
linear convolution coefficient has at most $M$ contributing pairs.  The factored real
product has bound $4M(B-1)^2$, while $2X\star Y$ has bound $2M(B-1)^2$.
Writing $[P]_j$ for coefficient~$j$ of a polynomial~$P$, the derivative products
satisfy $|[P_1]_j|,|[P_2]_j|\le M(B-1)^2$ and
$|[P_3]_j|\le 4M(B-1)^2$.  Consequently, the largest
conservative bound among the reconstructed nonlinear coefficients is
\begin{equation}
  |[2(P_3-P_1-P_2)]_j|
  \le 12M(B-1)^2
  \le 12\cdot 2^{25}(2^{16}-1)^2
  < \frac{p-1}{2}.
  \label{eq:reference-centered-headroom}
\end{equation}
The aligned constants are added after centered recovery, so they do not
consume this modular headroom.  Exactness here means exact coefficients for
the finite-precision inputs before the final significand truncation; it does
not make the entire iterated orbit exact over real numbers.

After centering, adjacent base-$2^{16}$ coefficients contribute to signed
base-$2^{32}$ limbs.  The fused B16 gather performs that conversion in the same
warp that constructs the carry metadata; release builds do not need to
materialize and reread four complete signed-limb arrays.  Debug builds may
materialize them so that the host and device checkpoints remain directly
comparable.

A carry chain is sequential if represented only as individual carry values.
It becomes parallel once a limb interval is represented as a transfer function
from its bounded incoming carry to its outgoing carry.  Function composition is
associative, so warp scans combine thread summaries, a block scan combines warp
summaries, and decoupled look-back resolves preceding block partitions.  The
method follows the parallel-prefix framework of Blelloch~\cite{Blelloch1990PrefixSums}
and the single-pass decoupled-look-back construction of Merrill and
Garland~\cite{merrill2016single}.

Concretely, for signed limb contribution $\ell_j$, incoming carry $\kappa_j$,
and base $Q=2^{32}$, the emitted digit and outgoing carry satisfy
\begin{equation}
  d_j=(\ell_j+\kappa_j)\bmod Q,
  \qquad
  \kappa_{j+1}=\frac{\ell_j+\kappa_j-d_j}{Q},
  \qquad 0\le d_j<Q.
  \label{eq:reference-signed-carry}
\end{equation}
Each limb therefore defines a small function
$\kappa_j\mapsto\kappa_{j+1}$.  Prefix composition finds every incoming
carry, after which the digits can be emitted independently, including when a
convolution coefficient or a linear contribution is negative.

The implementation packs the transfer functions for real, imaginary,
derivative-real, and derivative-imaginary streams into one 32-bit word in
FourWay mode.  A single composition instruction sequence therefore advances
four logically independent carry chains through the same scan topology.  After
the incoming carries are known, threads emit the final 32-bit digits in
parallel.

Normalization determines the sign, extracts the magnitude, finds the highest
nonzero digit, writes the fixed \code{GlobalNumUint32}-limb storage bucket, and
adjusts the exponent to preserve the retained bits' positional weights.
Subsequent arithmetic consumes only the most significant $D$ limbs of that
normalized bucket.  The next
iteration cannot begin from the inverse spectra alone: its alignment plan needs
this normalized sign, exponent, and significand.  Returning to
\code{HpSharkFloat} form is therefore a semantic dependency, not merely a
storage-format preference.

\subsection{Persistent Kernel and Host Lifecycle}
\label{subsec:ref-orbit-architecture}

Setup and execution have separate lifetimes.  The host selects a generated
storage family, chooses the active limb count, and constructs
\code{ReferencePreparedTables}.  Setup allocates one contiguous workspace,
assigns its spectrum and scratch regions, prepares plan descriptors from the
selected minimum through the selected maximum stage (stage~25 by default),
and launches the setup kernel to
generate their roots and twiddles.

\code{GpuOrbitSession} then owns the result descriptor and the CUDA stream for
one reference calculation.  Repeated \code{InvokeChunk} calls launch the
persistent recurrence for a requested number of iterations.  The production
callers limit each request to \code{MaxOutputIters}=1024 iterations, matching
the orbit-output capacity; \code{InvokeChunk} forwards the requested count.
The chunk limit bounds host transfer size and creates regular opportunities for cancellation and progress
handling without introducing a kernel launch for every recurrence.

Inside a launch, \code{HpSharkReferenceGpuLoop} obtains a cooperative
\code{grid\_group}.  Grid synchronization is required because a large NTT stage
can exchange data produced by different blocks.  A cooperative launch also
constrains the grid to blocks that can be resident together; launch selection
must therefore respect registers, dynamic shared memory, and the device's
cooperative occupancy.  The persistent organization amortizes launch overhead
and preserves the workspace across many sequential iterations, but it does not
remove the cost of the necessary barriers.

The generated CUDA translation units instantiate four families:
\code{HpSharkFloat\_\allowbreak Conversions}, \code{HpShark\allowbreak Reference},
\code{ReferenceGpu\allowbreak Loop}, and \code{SharkNTT\_\allowbreak Primitives}.  They provide
host/device conversion, lifecycle operations, the persistent recurrence, and
the Montgomery primitives shared with the host correctness model.  Explicit
instantiation keeps the supported precision families finite and makes their
workspace limits compile-time properties.

\subsection{Periodicity, Escape, and Resume Semantics}
\label{subsec:fused-ntt-periodicity}

When periodicity checking is enabled, the grid leader converts the current
high-precision orbit state to HDR values and evaluates the reduced-precision
criterion described in \cref{sec:hp-periodicity}.  Its derivative tracker is
also HDR, distinct from the FourWay high-precision derivative streams.
A continuing state allows the iteration plan to be built.  Escape, a positive periodicity result,
or an unrepresentable plan prevents the grid from entering another collective
phase.  This ordering is important: all threads must agree on whether they will
reach a cooperative barrier.

At loop entry, the ordinary reference path stores a low-precision sample of
the current $z_n$ and tests it before attempting the next recurrence.  The
leader increments the output count after a completed step, or when the
periodicity or escape check terminates on that sampled value.  The
next pass, if any, tests the newly published $z_{n+1}$.  The
Feature Finder instead returns the high-precision endpoint and derivative
state after its requested period.  Both clients use the same fused arithmetic
engine and differ in how they consume the results.

Ordinary reference construction starts the GPU state at $z=c$, with its HDR
derivative tracker initialized to one; the host supplies the preceding zero
orbit entry.  Feature Finder evaluation instead treats \code{startIter}=0 as
fresh initialization with $z_0=d_0=q_0=0$.  A positive \code{startIter} consumes
the supplied orbit and derivative values and continues from that checkpoint.  Abort and progress callbacks are
polled at bounded host-visible intervals without changing this contract.

\subsection{Checksum-Guided Host--Device Validation}
\label{sec:checksums}

Exact arithmetic makes bitwise stage comparison possible.  Debug builds can
record checksums after aligned packing, forward transforms, pointwise products,
inverse transforms, signed gathering, carry, and normalization.  The retained
\code{ReferenceReferenceOrbit2} CPU oracle executes the same plan and records
the corresponding host states.  \code{ChecksumsCheck} compares states by
purpose and call index, localizing a disagreement to a pipeline boundary
without retaining the obsolete v1 arithmetic operators.

The tests cover more than a forward/inverse round trip.  They compare orbit
outputs against MPIR, exercise exponent gaps and transform-length changes,
cover TwoWay and FourWay modes, verify checkpoint continuation, and run long
Newton sequences.  This breadth is required because an error in a late NTT
stage, a tail tile, or carry look-back can survive a short smoke test and become
visible only after many recurrence steps.

\subsection{Design Tradeoffs and Scope}
\label{subsec:fused-ntt-tradeoffs}

The fused v2 design is optimized for the regime in which high-precision
multiplication dominates reference construction.  Its main strengths are exact
field arithmetic, regular coefficient-level parallelism, transform reuse across
multiple streams, elimination of explicit bit reversal, and the absence of
normalized product temporaries.  Its costs are root and workspace preparation,
global barriers for large stages, modular-multiply latency, and occupancy
pressure from registers and shared memory.  A shorter operand may therefore be
better served by simpler arithmetic in another backend even though the NTT has
better asymptotic scaling.

Fusion also has a boundary.  Algebraic additions and derivative products can be
placed inside the shared transform and carry schedule, but independent orbit
iterations cannot be overlapped without changing the recurrence.  The inverse
transform, carry, and normalization cannot be moved to the end of the entire
orbit because the next iteration needs their result.  The implementation
concentrates on removing avoidable boundaries while retaining the mathematical
ones.

The application exposes one production selector for this engine.  The portable
menu labels it \emph{GPU-Accelerated Reference Orbit}, the Feature Finder labels
its backend \emph{NR Inner Loop: GPU}, and the CLI accepts \code{GPU}.
Reference-orbit rendering uses \code{GpuOrbitSession}; the Feature Finder calls
\code{EvaluateCriticalOrbitAnd\allowbreak Derivs\_\allowbreak GPU} through
\code{OrbitEndpoint\allowbreak Evaluator}.  The CUDA test harness retains
\code{Operator::\allowbreak ReferenceOrbit2} only as the established correctness-test
identity.  Production setup, execution, and shutdown use the unnumbered
\code{PrepareHpShark\allowbreak ReferenceTables},
\code{InvokeHpShark\allowbreak ReferenceKernel}, and
\code{ShutdownHpShark\allowbreak ReferenceKernel} interfaces.
